\section{Characteristic-Two Specializations and Kernel Consistency}\label{sec:binary}

This chapter develops the characteristic-two algebra required to
express and verify the masked correction of binary-field ring
switching in Boolean-kernel coordinates. The distinction from the
original KBFold protocol is essential: KBFold's smooth
multiplicative-domain word fold assumes odd characteristic, while
the trace-zero and Artin--Schreier identities here assume
characteristic two. We use the original Boolean-kernel basis
and normalized polynomial folds established in
Chapters~1, 4, and~5~\cite{MkhidaKBFold2026}.
The trace/Frobenius techniques should be compared with
Kochergin~\cite{Kochergin2026} and
Frigo--shelat~\cite{FrigoShelat2026};
we do not claim their trace-zero masking principle as new.

Every theorem is split into elementary lemmas with full proofs.
The principal new algebraic statements of this chapter are the
kernel-coordinate affine trace identity,
the independent anchored residual basis,
and a conditional two-challenge coefficient-consistency bound.
The final result is \emph{not} a zero-knowledge theorem for
the whole Ring Switching reduction or for the original KBFold PCS.

\subsection{Frobenius as a field and polynomial operation}\label{sec:binary:frobenius}
Throughout this chapter, $\F$ denotes a field of characteristic $2$. Unless explicitly assumed finite, the statements do not require that $\F$ have finitely many elements. For integers $j\geq0$ we write $q_j:=2^j$.

\begin{lemma}[Characteristic-two squaring]\label{lem:binary:square}
For all $a,b\in\F$, $(a+b)^2=a^2+b^2$.
\end{lemma}
\begin{proof}
Distributivity gives $(a+b)^2=a^2+ab+ba+b^2$. Since $\F$ is commutative and $2=0$, the middle terms sum to $2ab=0$.
\end{proof}

\begin{lemma}[Iterated Frobenius additivity]\label{lem:binary:frob-add}
For all $a,b\in\F$ and $j\geq0$,
\[
 (a+b)^{2^j}=a^{2^j}+b^{2^j}.
\]
\end{lemma}
\begin{proof}
Induct on $j$. At $j=0$ the identity is $a+b=a+b$. If it holds at $j$, square both sides and apply Lemma~\ref{lem:binary:square}:
\[
 (a+b)^{2^{j+1}}
 =\bigl(a^{2^j}+b^{2^j}\bigr)^2
 =a^{2^{j+1}}+b^{2^{j+1}}.
\]
\end{proof}

\begin{definition}[Frobenius on scalars]\label{def:binary:frob-scalar}
For $j\geq0$ define $\sigma_j:\F\to\F$ by $\sigma_j(a)=a^{2^j}$.
\end{definition}

\begin{lemma}[Frobenius endomorphism]\label{lem:binary:frob-ring}
For $j\geq0$, $\sigma_j$ preserves $0,1$, addition, and multiplication, and satisfies
$\sigma_{i+j}=\sigma_i\circ\sigma_j$ for all $i,j\geq0$.
\end{lemma}
\begin{proof}
The identities for $0,1$ are immediate. Additivity is Lemma~\ref{lem:binary:frob-add}. Multiplicativity follows from the exponent rule $(ab)^{2^j}=a^{2^j}b^{2^j}$ in a commutative field. Finally
$\sigma_i(\sigma_j(a))=(a^{2^j})^{2^i}=a^{2^{i+j}}=\sigma_{i+j}(a)$.
\end{proof}

\begin{lemma}[Frobenius is injective]\label{lem:binary:frob-injective}
For every $j\geq0$, $\sigma_j$ is injective.
\end{lemma}
\begin{proof}
If $\sigma_j(a)=\sigma_j(b)$, then by Lemma~\ref{lem:binary:frob-add},
$\sigma_j(a+b)=0$. Because a field has no zero divisors, $x^{2^j}=0$ implies $x=0$; therefore $a+b=0$, or $a=b$ in characteristic two.
\end{proof}

\begin{corollary}[Frobenius permutes a finite field]\label{cor:binary:frob-bijective}
If $\F$ is finite, every $\sigma_j$ is bijective, and sends a uniform $\F$-valued variable to another uniform variable.
\end{corollary}
\begin{proof}
An injective map from a finite set to itself is surjective, proving bijectivity. A bijection preserves the uniform distribution because every output has exactly one preimage.
\end{proof}

\begin{definition}[Coefficient Frobenius of a polynomial]\label{def:binary:poly-frob}
For $P(X)=\sum_{e=0}^{d}p_eX^e\in\F[X]$, define
\[
 P^{\sigma_j}(X):=\sum_{e=0}^{d}p_e^{2^j}X^e.
\]
This notation acts on coefficients and does not replace $X$ by $X^{2^j}$.
\end{definition}

\begin{lemma}[Polynomial Frobenius identity]\label{lem:binary:poly-frob}
For every $P\in\F[X]$ and $j\geq0$,
\[
 \boxed{P(X)^{2^j}=P^{\sigma_j}(X^{2^j}).}
\]
\end{lemma}
\begin{proof}
Write $P=\sum_ep_eX^e$. Repeated application of Lemma~\ref{lem:binary:square} in the polynomial ring $\F[X]$ gives
$P(X)^{2^j}=\sum_e(p_eX^e)^{2^j}$. Exponent rules yield
$\sum_ep_e^{2^j}X^{e2^j}$, which equals
$\sum_ep_e^{2^j}(X^{2^j})^e=P^{\sigma_j}(X^{2^j})$ by definition.
\end{proof}

\begin{lemma}[Coefficient Frobenius and products]\label{lem:binary:poly-frob-mult}
For $P,Q\in\F[X]$ and $j\geq0$,
\[
 (P+Q)^{\sigma_j}=P^{\sigma_j}+Q^{\sigma_j},\qquad
 (PQ)^{\sigma_j}=P^{\sigma_j}Q^{\sigma_j}.
\]
\end{lemma}
\begin{proof}
The first identity follows coefficientwise from Lemma~\ref{lem:binary:frob-add}. The coefficient of $X^e$ in $PQ$ is $\sum_{a+b=e}p_aq_b$; applying $\sigma_j$ and using its additivity and multiplicativity gives
$\sum_{a+b=e}p_a^{2^j}q_b^{2^j}$, which is the coefficient of $X^e$ in $P^{\sigma_j}Q^{\sigma_j}$.
\end{proof}

\begin{remark}[Formal distinction]
The identity in Lemma~\ref{lem:binary:poly-frob} is an equality of ordinary polynomials. The operation $P\mapsto P^{\sigma_j}$ and the substitution $P\mapsto P(X^{2^j})$ have different types in Lean and must not be conflated.
\end{remark}

\subsection{Frobenius stability and the Boolean kernel partition}\label{sec:binary:kernel}

Fix $\ell\geq1$ and put $n:=2^\ell$. Let $I_\ell=\{0,1\}^\ell$, and let
$K_b(X)=\prod_{i=0}^{\ell-1}(X^{2^i}+b_i)$ denote the Boolean kernel from Definition~\ref{def:intro:kernel}.

\begin{lemma}[Binary coefficients of a kernel]\label{lem:binary:binary-coeffs}
For every $b\in I_\ell$, the polynomial $K_b$ has all coefficients in
the prime subfield $\F_2\subseteq\F$.
\end{lemma}
\begin{proof}
Every factor $X^{2^i}+b_i$ has coefficients in $\F_2$ because
$b_i\in\{0,1\}$. A finite product of polynomials with coefficients
in a subfield again has coefficients in that subfield, by the
finite-sum formula for multiplication coefficients.
\end{proof}

\begin{lemma}[Kernel coefficient Frobenius invariance]\label{lem:binary:kernel-coeff-frob}
For $b\in I_\ell$ and $j\geq0$,
\[
 K_b^{\sigma_j}(X)=K_b(X).
\]
\end{lemma}
\begin{proof}
By Lemma~\ref{lem:binary:binary-coeffs}, each coefficient is
$0$ or $1$. Both are fixed by $\sigma_j$, so the coefficientwise
Frobenius operation leaves the polynomial unchanged.
\end{proof}

\begin{theorem}[Kernel Frobenius substitution]\label{thm:binary:kernel-square}
For every $b\in I_\ell$ and $j\geq0$,
\[
 \boxed{K_b(X)^{2^j}=K_b(X^{2^j}).}
\]
\end{theorem}
\begin{proof}
Lemma~\ref{lem:binary:poly-frob} gives
$K_b(X)^{2^j}=K_b^{\sigma_j}(X^{2^j})$.
By Lemma~\ref{lem:binary:kernel-coeff-frob} the latter equals
$K_b(X^{2^j})$.
\end{proof}

\begin{definition}[Kernel coefficient Frobenius]\label{def:binary:table-frob}
For a coefficient table $c:I_\ell\to\F$ and $j\geq0$, define
$c^{\sigma_j}(b):=c(b)^{2^j}$ and
$U_c(X):=\sum_{b\in I_\ell}c(b)K_b(X)$.
\end{definition}

\begin{theorem}[Frobenius covariance of kernel encoding]\label{thm:binary:encoding-frob}
For every coefficient table $c$ and $j\geq0$,
\[
 \boxed{U_c(X)^{2^j}=U_{c^{\sigma_j}}(X^{2^j}).}
\]
\end{theorem}
\begin{proof}
Repeatedly apply Lemma~\ref{lem:binary:square} to the finite sum
defining $U_c$:
\[
 U_c(X)^{2^j}
 =\sum_b c(b)^{2^j}K_b(X)^{2^j}.
\]
Theorem~\ref{thm:binary:kernel-square} turns each kernel factor
into $K_b(X^{2^j})$, giving
$\sum_b c^{\sigma_j}(b)K_b(X^{2^j})$,
which is $U_{c^{\sigma_j}}(X^{2^j})$.
\end{proof}

\begin{lemma}[Kernel partition sum]\label{lem:binary:partition}
\[
 \boxed{\sum_{b\in I_\ell}K_b(X)=1.}
\]
\end{lemma}
\begin{proof}
Distribute the sum over all bit strings:
\[
 \sum_{b\in I_\ell}\prod_{i=0}^{\ell-1}(X^{2^i}+b_i)
 =\prod_{i=0}^{\ell-1}
    \bigl((X^{2^i}+0)+(X^{2^i}+1)\bigr).
\]
Inside each parenthesis, $X^{2^i}+X^{2^i}=0$ in
characteristic two, leaving the scalar $1$. The product of
$\ell$ copies of $1$ is $1$.
\end{proof}

\begin{corollary}[All-one kernel coordinates of the constant polynomial]\label{cor:binary:one-coeffs}
Let $\mathbf1:I_\ell\to\F$ be the constant table with value $1$.
Then $U_{\mathbf1}(X)=1$. Moreover the coordinates of $1$
in the kernel basis are uniquely all equal to $1$.
\end{corollary}
\begin{proof}
By definition $U_{\mathbf1}=\sum_bK_b=1$,
using Lemma~\ref{lem:binary:partition}. Uniqueness follows
from the kernel basis theorem, Proposition~\ref{prop:intro:basis}.
\end{proof}

\begin{lemma}[Two complementary lowest-bit kernels]\label{lem:binary:complementary}
For every $b'\in I_{\ell-1}$,
\[
 K_{(0,b')}(X)+K_{(1,b')}(X)=K_{b'}^{(\ell-1)}(X^2).
\]
\end{lemma}
\begin{proof}
By Lemma~\ref{lem:intro:recursion},
$K_{(0,b')}=XK_{b'}^{(\ell-1)}(X^2)$ and
$K_{(1,b')}=(X+1)K_{b'}^{(\ell-1)}(X^2)$.
Add and factor out the common polynomial. The scalar
$X+(X+1)$ equals $1$ because $X+X=0$.
\end{proof}

\begin{remark}[One important characteristic restriction]
Lemma~\ref{lem:binary:partition} is false over a general
odd-characteristic field. The exact all-characteristic identity
is Proposition~\ref{prop:kbalgebra:partition}. We will use
$\sum_bK_b=1$ to identify the degeneracy of an \emph{unanchored}
residual polynomial, so the characteristic assumption cannot
be dropped.
\end{remark}

\subsection{Multilinearization and the exact kernel-fold functional}\label{sec:binary:multilinearization}
We retain $\ell\geq1$, $n=2^\ell$, and a characteristic-two field $\F$.
For $0\leq e<n$ write its unique binary expansion as
$e=\sum_{i=0}^{\ell-1}e_i2^i$ with $e_i\in\{0,1\}$.

\begin{definition}[Binary-exponent multilinearization]\label{def:binary:ml}
Let $\mathscr{ML}_\ell$ be the $\F$-vector space of
polynomials multilinear in $Z_0,\ldots,Z_{\ell-1}$.
Define the $\F$-linear map
\[
 \mathscr M_\ell:\F[X]_{<n}\longrightarrow\mathscr{ML}_\ell,
 \qquad
 \mathscr M_\ell\Bigl(\sum_{e=0}^{n-1}a_eX^e\Bigr)
 :=\sum_{e=0}^{n-1}a_e\prod_{i=0}^{\ell-1}Z_i^{e_i}.
\]
\end{definition}

\begin{lemma}[Multilinearization is a vector-space isomorphism]\label{lem:binary:ml-iso}
The map $\mathscr M_\ell$ is an $\F$-linear bijection.
\end{lemma}
\begin{proof}
Linearity is part of its definition. The monomials
$X^e$, $0\leq e<n$, form a basis of $\F[X]_{<n}$.
The monomials $\prod_iZ_i^{e_i}$ indexed by the binary digits of
$0\leq e<n$ form a basis of $\mathscr{ML}_\ell$.
The binary expansion gives a bijection between these two bases.
The map sends each element of the first basis to the corresponding
element of the second, hence is an isomorphism.
\end{proof}

\begin{lemma}[Explicit subset expansion of a kernel]\label{lem:binary:kernel-expansion}
For $b\in I_\ell$,
\[
 K_b(X)=\sum_{S\subseteq\{0,\ldots,\ell-1\}}
   \left(\prod_{i\notin S}b_i\right)
   X^{\sum_{i\in S}2^i}.
\]
\end{lemma}
\begin{proof}
In the product $K_b=\prod_i(X^{2^i}+b_i)$,
a term is obtained by selecting $X^{2^i}$ from a subset $S$
of factors and $b_i$ from the complementary factors.
The resulting term has scalar coefficient $\prod_{i\notin S}b_i$
and monomial factor $X^{\sum_{i\in S}2^i}$.
Every subset occurs exactly once in the full distributive
expansion, establishing the equality.
\end{proof}

\begin{lemma}[Multilinearization of every kernel basis element]\label{lem:binary:ml-kernel}
For every $b\in I_\ell$,
\[
 \boxed{\mathscr M_\ell(K_b)(Z)=\prod_{i=0}^{\ell-1}(Z_i+b_i).}
\]
\end{lemma}
\begin{proof}
By Lemma~\ref{lem:binary:kernel-expansion} and linearity,
\[
 \mathscr M_\ell(K_b)=\sum_{S}
   \left(\prod_{i\notin S}b_i\right)
   \mathscr M_\ell\left(X^{\sum_{i\in S}2^i}\right).
\]
The binary digits of $\sum_{i\in S}2^i$ equal $1$ precisely for
$i\in S$, so the last factor is $\prod_{i\in S}Z_i$.
Therefore the sum equals
$\sum_S(\prod_{i\notin S}b_i)(\prod_{i\in S}Z_i)$.
The distributive expansion of $\prod_i(Z_i+b_i)$ is the same sum.
This argument uses only linearity of $\mathscr M_\ell$,
not multiplicativity.
\end{proof}

\begin{definition}[Characteristic-two equality weights]\label{def:binary:lambda}
For $b\in I_\ell$ and $\theta=(\theta_0,\ldots,\theta_{\ell-1})\in\F^\ell$,
define
\[
 \lambda_b(\theta):=\prod_{i=0}^{\ell-1}(1+\theta_i+b_i).
\]
For $P=\sum_bc_bK_b$ with its unique kernel coefficient table,
define
\[
 D_P(\theta):=\sum_bc_b\lambda_b(\theta).
\]
\end{definition}

\begin{lemma}[Equality weights agree]\label{lem:binary:eq-weights}
For all $b\in I_\ell$ and $\theta\in\F^\ell$,
$\lambda_b(\theta)=\operatorname{eq}_\ell(b,\theta)$.
\end{lemma}
\begin{proof}
For each coordinate $i$, if $b_i=0$ then
$1+\theta_i+b_i=1+\theta_i=1-\theta_i$.
If $b_i=1$ then $1+\theta_i+b_i=\theta_i$,
because $1+1=0$. These are precisely the two cases of
$(1-b_i)(1-\theta_i)+b_i\theta_i$.
Multiply the coordinatewise identities.
\end{proof}

\begin{lemma}[Boolean weights partition unity]\label{lem:binary:lambda-partition}
For every $\theta\in\F^\ell$,
$\sum_{b\in I_\ell}\lambda_b(\theta)=1$.
\end{lemma}
\begin{proof}
Factor the sum over bit strings into a product of two-term sums:
\[
 \sum_b\lambda_b(\theta)
 =\prod_{i=0}^{\ell-1}
  \bigl((1+\theta_i+0)+(1+\theta_i+1)\bigr)
 =\prod_i\bigl((1+\theta_i)+\theta_i\bigr)=1.
\]
\end{proof}

\begin{theorem}[Complete kernel fold equals multilinearized evaluation]\label{thm:binary:fold-ml}
For every $P\in\F[X]_{<n}$ and $\theta\in\F^\ell$,
\[
 \boxed{
 \mathsf{PFold}^{[\ell]}_\theta(P)
  =D_P(\theta)
  =\mathscr M_\ell(P)(1+\theta_0,\ldots,1+\theta_{\ell-1}).
 }
\]
\end{theorem}
\begin{proof}
By Proposition~\ref{prop:intro:basis}, uniquely write
$P=\sum_bc_bK_b$. The fold dictionary,
Corollary~\ref{cor:kbff:full-evaluation},
identifies its complete fold with
$\sum_bc_b\operatorname{eq}_\ell(b,\theta)$,
which equals $D_P(\theta)$ by Lemma~\ref{lem:binary:eq-weights}.
For the second equality, apply the linear map $\mathscr M_\ell$
to the kernel expansion of $P$ and evaluate:
\[
 \mathscr M_\ell(P)(1+\theta)
 =\sum_bc_b\,\mathscr M_\ell(K_b)(1+\theta).
\]
Lemma~\ref{lem:binary:ml-kernel} evaluates each summand as
$\prod_i(1+\theta_i+b_i)=\lambda_b(\theta)$.
Thus the sum is precisely $D_P(\theta)$.
\end{proof}

\begin{definition}[Partial multilinearization]\label{def:binary:partial-ml}
For $0\leq j\leq\ell$, define
\[
 \mathscr M_{\ell,j}(P)(Z_0,\ldots,Z_{j-1},Y)
 :=\sum_{e=0}^{n-1}a_e
   \left(\prod_{i=0}^{j-1}Z_i^{e_i}\right)
   Y^{\lfloor e/2^j\rfloor}
\]
for $P(X)=\sum_{e=0}^{n-1}a_eX^e$.
\end{definition}

\begin{lemma}[Partial folding formula]\label{lem:binary:partial-fold}
For $0\leq j\leq\ell$ and $r\in\F^j$,
\[
 \boxed{
 \mathsf{PFold}^{[j]}_r(P)(Y)
 =\mathscr M_{\ell,j}(P)(1+r_0,\ldots,1+r_{j-1},Y).
 }
\]
\end{lemma}
\begin{proof}
The explicit fold formula of
Theorem~\ref{thm:kbff:fold-coefficients}
gives a summand
$a_e\bigl(\prod_{i<j}A_{e_i}(r_i)\bigr)Y^{\lfloor e/2^j\rfloor}$.
In characteristic two, $A_0(r_i)=2r_i-1=1$ and
$A_1(r_i)=1-r_i=1+r_i$.
Thus $\prod_{i<j}A_{e_i}(r_i)=\prod_{i<j}(1+r_i)^{e_i}$.
Substitution produces exactly the evaluation of
Definition~\ref{def:binary:partial-ml}.
\end{proof}

\begin{lemma}[Frobenius covariance of equality weights]\label{lem:binary:lambda-frob}
For $b\in I_\ell$, $\theta\in\F^\ell$, and $j\geq0$,
\[
 \lambda_b(\theta)^{2^j}
 =\lambda_b(\theta_0^{2^j},\ldots,\theta_{\ell-1}^{2^j}).
\]
\end{lemma}
\begin{proof}
Apply $\sigma_j$ to the product
$\prod_i(1+\theta_i+b_i)$, using the ring-homomorphism
property of Lemma~\ref{lem:binary:frob-ring}.
Since $1$ and $b_i\in\F_2$ are fixed by $\sigma_j$,
each factor becomes $1+\theta_i^{2^j}+b_i$.
Their product is the equality weight on the right.
\end{proof}

\begin{corollary}[Frobenius covariance of kernel-fold values]\label{cor:binary:fold-frob}
For $c:I_\ell\to\F$, $\theta\in\F^\ell$, and $j\geq0$,
\[
 \boxed{D_{U_c}(\theta)^{2^j}
 =D_{U_{c^{\sigma_j}}}(\theta^{\sigma_j}).}
\]
\end{corollary}
\begin{proof}
Raise the finite sum
$D_{U_c}(\theta)=\sum_bc(b)\lambda_b(\theta)$
to its $2^j$-th power.
By Lemma~\ref{lem:binary:frob-add} and multiplicativity,
it becomes $\sum_bc(b)^{2^j}\lambda_b(\theta)^{2^j}$.
Apply Lemma~\ref{lem:binary:lambda-frob}
and Definition~\ref{def:binary:table-frob}.
\end{proof}

\begin{remark}[Nonmultiplicativity warning]
The multilinearization map is \emph{not} a ring homomorphism,
even where both sides happen to be defined.
For $\ell\geq2$,
$\mathscr M_\ell(X^2)=Z_1$ while
$\mathscr M_\ell(X)\mathscr M_\ell(X)=Z_0^2$.
The polynomials $Z_1$ and $Z_0^2$ are distinct.
All foregoing kernel identities were proved from expansions
and the linearity of $\mathscr M_\ell$.
\end{remark}

\subsection{Finite binary extensions, trace, and Artin--Schreier masks}\label{sec:binary:trace}

For the rest of the chapter, specialize to
$L=\F_{2^n}$ with $n=2^\ell$ and $\ell\geq1$.
The statements that mention uniformity use the finite uniform
distribution from Definition~\ref{def:security:distribution}.
We write $q:=|L|=2^n$.

\begin{lemma}[Frobenius has period $n$]\label{lem:binary:frob-period}
For every $a\in L$, $a^{2^n}=a$.
\end{lemma}
\begin{proof}
If $a=0$, the identity is immediate. Suppose $a\neq0$.
Multiplication by $a$ is a bijection of $L^\times$,
whose cardinality is $q-1$. Define
$P:=\prod_{x\in L^\times}x$, a nonzero element of $L$.
Reindexing the product under the bijection $x\mapsto ax$ gives
\[
 P=\prod_{x\in L^\times}(ax)
  =a^{q-1}\prod_{x\in L^\times}x=a^{q-1}P.
\]
Cancel the nonzero element $P$, obtaining $a^{q-1}=1$.
Multiplying by $a$ yields $a^q=a$, with $q=2^n$.
\end{proof}

\begin{definition}[Absolute trace]\label{def:binary:trace}
For $x\in L$ define
\[
 \Tr(x):=\sum_{j=0}^{n-1}x^{2^j}.
\]
Initially the sum is taken in $L$; the next lemma proves its
value lies in the prime subfield $\F_2$.
\end{definition}

\begin{lemma}[Trace is additive and Frobenius-invariant]\label{lem:binary:trace-add}
For every $x,y\in L$,
\[
 \Tr(x+y)=\Tr(x)+\Tr(y),\qquad
 \Tr(x^2)=\Tr(x).
\]
In particular, $\Tr$ is $\F_2$-linear.
\end{lemma}
\begin{proof}
The first equality follows by applying
Lemma~\ref{lem:binary:frob-add} to each term of the finite sum.
For the second, write
\[
 \Tr(x^2)=x^2+x^{2^2}+\cdots+x^{2^n}.
\]
By Lemma~\ref{lem:binary:frob-period}, the last term is $x$.
Thus the sum is a cyclic permutation of the terms in $\Tr(x)$.
Linearity for the scalar $0$ is immediate, and for scalar $1$
it is the identity $\Tr(1x)=\Tr(x)$.
\end{proof}

\begin{lemma}[Trace takes binary values]\label{lem:binary:trace-binary}
For every $x\in L$, $\Tr(x)\in\F_2=\{0,1\}$.
\end{lemma}
\begin{proof}
By Lemma~\ref{lem:binary:trace-add},
$\Tr(x)^2=\Tr(x^2)=\Tr(x)$.
Hence $\Tr(x)$ is a root of $T^2+T=T(T+1)$.
A field has no zero divisors, so the only possibilities are
$T=0$ and $T=1$, both in $\F_2$.
\end{proof}

\begin{lemma}[Trace is surjective]\label{lem:binary:trace-surj}
The linear map $\Tr:L\to\F_2$ is nonzero and therefore surjective.
\end{lemma}
\begin{proof}
Consider the polynomial
$H(X)=X+X^2+\cdots+X^{2^{n-1}}\in L[X]$.
It is nonzero: its coefficient of $X$ is $1$, and all higher
exponents are distinct. Its degree $2^{n-1}$ is strictly less than
$q=2^n$. If $H(x)=0$ for all $x\in L$, this nonzero polynomial
would have $q$ distinct roots, contradicting
Lemma~\ref{lem:history:rootcount}. Thus some $x$ satisfies
$\Tr(x)=H(x)\neq0$. By Lemma~\ref{lem:binary:trace-binary},
that value equals $1$. Therefore the image of $\Tr$ contains
both $0$ and $1$.
\end{proof}

\begin{corollary}[Sizes of the two trace fibers]\label{cor:binary:trace-fibers}
The sets $\ker\Tr$ and $\{x\in L:\Tr(x)=1\}$ each have
$q/2$ elements.
\end{corollary}
\begin{proof}
Choose $a\in L$ with $\Tr(a)=1$ by
Lemma~\ref{lem:binary:trace-surj}. Translation $x\mapsto x+a$
is a bijection and, by trace additivity, sends elements of trace
zero exactly to elements of trace one and vice versa.
The two fibers partition $L$, and hence have equal cardinality
$q/2$.
\end{proof}

\begin{definition}[Artin--Schreier operator]\label{def:binary:artin-schreier}
Define $\wp:L\to L$ by
\[
 \wp(t):=t^2+t.
\]
\end{definition}

\begin{lemma}[Artin--Schreier is binary-linear]\label{lem:binary:as-linear}
For all $s,t\in L$, $\wp(s+t)=\wp(s)+\wp(t)$.
\end{lemma}
\begin{proof}
Expand $(s+t)^2+s+t$ and apply
Lemma~\ref{lem:binary:square}. The result is
$(s^2+s)+(t^2+t)$.
\end{proof}

\begin{lemma}[Artin--Schreier kernel]\label{lem:binary:as-kernel}
$\ker\wp=\{0,1\}$.
\end{lemma}
\begin{proof}
$\wp(t)=t^2+t=t(t+1)$. This product is zero exactly when
$t=0$ or $t+1=0$, the latter being $t=1$ in characteristic two.
Both values satisfy the equation.
\end{proof}

\begin{lemma}[Every Artin--Schreier value has trace zero]\label{lem:binary:as-trace}
For every $t\in L$, $\Tr(\wp(t))=0$.
\end{lemma}
\begin{proof}
By trace additivity and Frobenius invariance,
$\Tr(t^2+t)=\Tr(t^2)+\Tr(t)=\Tr(t)+\Tr(t)=0$.
\end{proof}

\begin{theorem}[Artin--Schreier characterizes trace zero]\label{thm:binary:as-image}
\[
 \boxed{\operatorname{im}\wp=\ker\Tr.}
\]
Every element of $\ker\Tr$ has exactly two preimages under $\wp$.
\end{theorem}
\begin{proof}
Lemma~\ref{lem:binary:as-trace} gives
$\operatorname{im}\wp\subseteq\ker\Tr$.
By Lemma~\ref{lem:binary:as-linear},
the nonempty fiber above any $y=\wp(t)$ is the coset
$t+\ker\wp$: indeed $\wp(s)=\wp(t)$ is equivalent to
$\wp(s+t)=0$ and hence $s+t\in\ker\wp$.
By Lemma~\ref{lem:binary:as-kernel}, every nonempty fiber
contains exactly two elements. Thus
$|\operatorname{im}\wp|=q/2$.
Corollary~\ref{cor:binary:trace-fibers} shows
$|\ker\Tr|=q/2$. The inclusion between equal-sized finite
sets is an equality.
\end{proof}

\begin{corollary}[Uniform trace-zero masking]\label{cor:binary:as-uniform}
If $T$ is uniform on $L$, then $\wp(T)$ is uniform on
$\ker\Tr$.
\end{corollary}
\begin{proof}
By Theorem~\ref{thm:binary:as-image}, every trace-zero
element has exactly two preimages. Therefore its probability
as the output $\wp(T)$ is $2/q=1/(q/2)$, independent of
the chosen trace-zero element.
\end{proof}

\begin{theorem}[Perfect hiding of an individual trace-zero shift]\label{thm:binary:single-mask}
Fix $v\in\ker\Tr$ and let $T$ be uniform on $L$.
Then
\[
 U:=v+\wp(T)
\]
is uniform on $\ker\Tr$, independently of $v$.
\end{theorem}
\begin{proof}
Corollary~\ref{cor:binary:as-uniform} makes
$\wp(T)$ uniform on the additive subgroup $\ker\Tr$.
Since $v$ lies in that subgroup, translation
$y\mapsto v+y$ is a bijection of $\ker\Tr$,
so $U$ is uniform there. Its distribution is therefore
independent of the fixed value $v$.
\end{proof}

\begin{remark}[No joint-ZK conclusion]
Theorem~\ref{thm:binary:single-mask} concerns one public
field element. If the same $T$ is subsequently used in
another published equation, the joint transcript can have
correlations. The complete verification protocol must
analyze those correlations, as illustrated by the
reused-mask attack of Theorem~\ref{thm:kbff:one-query-attack}.
\end{remark}

\subsection{Trace-dual coordinates and affine verification of kernel corrections}\label{sec:binary:coordinates}
All vector indices in this subsection belong to $I_\ell$ and all
field elements belong to $L=\F_{2^n}$, $n=2^\ell$.
The characteristic-two coordinate construction is related
to binary-field ring switching~\cite{DiamondPosen2024}
and to its masked versions~\cite{Kochergin2026,FrigoShelat2026}.

\begin{lemma}[Existence of an element of full binary degree]\label{lem:binary:full-degree-element}
There exists $\alpha\in L=\F_{2^n}$ such that
$[\F_2(\alpha):\F_2]=n$.
\end{lemma}
\begin{proof}
For any $a\in L$, the tower law applied to
$\F_2\subseteq\F_2(a)\subseteq L$
shows that $d:=[\F_2(a):\F_2]$ divides $n$.
If every $a$ had $d<n$, then, because $n=2^\ell$,
every such $d$ would divide $n/2$.
The field $\F_2(a)$ has exactly $2^d$ elements.
Repeating the product-permutation proof of
Lemma~\ref{lem:binary:frob-period} in this
finite field gives $a^{2^d}=a$.
Iterating Frobenius then yields
$a^{2^{n/2}}=a$ for every $a\in L$.
This would make all $2^n$ elements of $L$ roots of
$X^{2^{n/2}}+X$, a nonzero polynomial of degree
$2^{n/2}<2^n$. It contradicts the root-count lemma,
Lemma~\ref{lem:history:rootcount}.
Therefore some $\alpha$ has degree $n$.
\end{proof}

\begin{definition}[Evaluation basis]\label{def:binary:eval-basis}
Choose $\alpha\in L$ of degree $n$ over $\F_2$,
whose existence is guaranteed by
Lemma~\ref{lem:binary:full-degree-element}, and set
\[
 \beta_u:=K_u(\alpha),\qquad u\in I_\ell.
\]
\end{definition}

\begin{lemma}[Kernel values form a binary field basis]\label{lem:binary:eval-basis}
The family $(\beta_u)_{u\in I_\ell}$ is an $\F_2$-basis of $L$.
\end{lemma}
\begin{proof}
The polynomials $(K_u)_u$ form an $\F_2$-basis of
$\F_2[X]_{<n}$ by Proposition~\ref{prop:intro:basis}.
The evaluation map $P\mapsto P(\alpha)$ is injective on
that space: if a nonzero polynomial of degree $<n$
vanished at $\alpha$, it would contradict the minimality
of the degree-$n$ minimal polynomial.
Both domain and codomain have dimension $n$ over $\F_2$.
Thus the map is an isomorphism, and the images of the
kernel basis form a basis of $L$.
\end{proof}

\begin{definition}[Binary coordinate maps]\label{def:binary:coords}
For each $u\in I_\ell$, define the unique
$\F_2$-linear coordinate map $d_u:L\to\F_2$ by
\[
 x=\sum_{u\in I_\ell}\beta_u d_u(x)
 \quad\text{for every }x\in L.
\]
\end{definition}

\begin{lemma}[Coordinate extraction]\label{lem:binary:extract}
For $u,v\in I_\ell$, $d_u(\beta_v)=1$ if $u=v$ and $0$
otherwise; each $d_u$ is additive.
\end{lemma}
\begin{proof}
The kernel-basis expansion of $\beta_v$ has coefficient $1$
at index $v$ and $0$ elsewhere; uniqueness of coordinates
gives the first assertion. If $x=\sum_v\beta_v a_v$ and
$y=\sum_v\beta_v b_v$, with $a_v,b_v\in\F_2$, then
$x+y=\sum_v\beta_v(a_v+b_v)$. Uniqueness gives
$d_u(x+y)=a_u+b_u=d_u(x)+d_u(y)$.
\end{proof}

\begin{lemma}[Trace pairing is nondegenerate]\label{lem:binary:trace-pairing}
For every nonzero $a\in L$ there exists $x\in L$ with
$\Tr(ax)=1$.
\end{lemma}
\begin{proof}
Choose $y\in L$ with $\Tr(y)=1$, using
Lemma~\ref{lem:binary:trace-surj}. As $a\neq0$,
$a^{-1}$ exists. Put $x=a^{-1}y$; then
$\Tr(ax)=\Tr(y)=1$.
\end{proof}

\begin{theorem}[Existence and uniqueness of trace-dual coordinates]\label{thm:binary:dual-basis}
There is a unique family $(\eta_u)_{u\in I_\ell}$ in $L$ such that
\[
 \boxed{\Tr(\eta_u\beta_v)=\delta_{uv}.}
\]
For every $x\in L$,
\[
 \boxed{d_u(x)=\Tr(\eta_u x).}
\]
\end{theorem}
\begin{proof}
Define the $\F_2$-linear map
$\Phi:L\to\operatorname{Hom}_{\F_2}(L,\F_2)$ by
$\Phi(a)(x)=\Tr(ax)$. Linearity follows from
Lemma~\ref{lem:binary:trace-add}.
Lemma~\ref{lem:binary:trace-pairing} proves that $\Phi$
is injective. Its domain has dimension $n$ over $\F_2$;
its codomain also has dimension $n$, since a linear functional
is uniquely determined by its values on the $n$ basis elements
$\beta_v$. Hence $\Phi$ is bijective.
For each coordinate functional $d_u$, choose the unique
$\eta_u=\Phi^{-1}(d_u)$. Then
$\Tr(\eta_u\beta_v)=d_u(\beta_v)=\delta_{uv}$
by Lemma~\ref{lem:binary:extract}, and
$\Tr(\eta_u x)=d_u(x)$ for all $x$ by construction.
Uniqueness of each $\eta_u$ follows from injectivity
of $\Phi$.
\end{proof}

\begin{definition}[Binary transpose and kernel correction]\label{def:binary:transpose}
For $p=(p_v)_{v\in I_\ell}\in L^{I_\ell}$ define
\[
 c^p_u:=\sum_{v\in I_\ell}\beta_vd_u(p_v),\qquad
 C_p(X):=\sum_{u\in I_\ell}c^p_uK_u(X).
\]
For $\theta\in L^\ell$ define
\[
 \Delta_p(\theta):=\sum_{u\in I_\ell}
       \lambda_u(\theta)c^p_u
       =D_{C_p}(\theta).
\]
\end{definition}

\begin{lemma}[Trace expression for the correction]\label{lem:binary:delta-trace}
For every $p\in L^{I_\ell}$ and $\theta\in L^\ell$,
\[
 \boxed{\Delta_p(\theta)=
  \sum_{u,v\in I_\ell}
   \lambda_u(\theta)\beta_v\Tr(\eta_u p_v).}
\]
\end{lemma}
\begin{proof}
Expand $\Delta_p(\theta)$ using
Definition~\ref{def:binary:transpose}:
\[
 \Delta_p(\theta)=
 \sum_u\lambda_u(\theta)\sum_v\beta_vd_u(p_v).
\]
Interchange the two finite sums and replace
$d_u(p_v)$ by $\Tr(\eta_up_v)$ using
Theorem~\ref{thm:binary:dual-basis}.
\end{proof}

\begin{definition}[Public affine coefficients]\label{def:binary:affine-H}
For $\theta\in L^\ell$ and $u,v,k\in I_\ell$ put
\[
 H_{v,k}(\theta):=\sum_{u\in I_\ell}
       d_k\bigl(\beta_v\lambda_u(\theta)\bigr)\eta_u.
\]
For a candidate coefficient table $c:I_\ell\to L$ define
\[
 D_c(\theta):=\sum_u\lambda_u(\theta)c_u
\]
and
\[
 V_k(p,c,\theta):=\eta_kD_c(\theta)
      +\sum_{v\in I_\ell}H_{v,k}(\theta)p_v.
\]
The coefficients $H_{v,k}(\theta)$ depend only on public
parameters and $\theta$, not on the hidden pads $p_v$.
\end{definition}

\begin{lemma}[Affine transpose trace identity]\label{lem:binary:affine-trace}
For every $p,\theta,k$,
\[
 \Tr\left(\sum_v H_{v,k}(\theta)p_v\right)
   =d_k\bigl(\Delta_p(\theta)\bigr).
\]
\end{lemma}
\begin{proof}
Expand the definition of $H$ and use
$\F_2$-linearity of trace:
\[
 \Tr\left(\sum_vH_{v,k}p_v\right)
 =\sum_{u,v}
    d_k(\beta_v\lambda_u(\theta))\,\Tr(\eta_up_v).
\]
Each factor $\Tr(\eta_up_v)$ lies in $\F_2$.
Using additivity and $\F_2$-linearity of $d_k$,
the last sum equals
\[
 d_k\left(\sum_{u,v}
   \beta_v\lambda_u(\theta)\Tr(\eta_up_v)\right).
\]
The field element inside $d_k$ equals
$\Delta_p(\theta)$ by Lemma~\ref{lem:binary:delta-trace}.
\end{proof}

\begin{lemma}[Trace coordinates of affine residuals]\label{lem:binary:V-trace}
For all $p,c,\theta$ and $k$,
\[
 \boxed{\Tr(V_k(p,c,\theta))
      =d_k\bigl(D_c(\theta)+\Delta_p(\theta)\bigr).}
\]
\end{lemma}
\begin{proof}
By the definition of $V_k$ and trace additivity,
\[
 \Tr(V_k)=\Tr(\eta_k D_c(\theta))
  +\Tr\left(\sum_vH_{v,k}p_v\right).
\]
The first term equals $d_k(D_c(\theta))$ by
Theorem~\ref{thm:binary:dual-basis}; the second equals
$d_k(\Delta_p(\theta))$ by
Lemma~\ref{lem:binary:affine-trace}.
Additivity of $d_k$ combines them into the claimed formula.
\end{proof}

\begin{theorem}[Exact affine trace consistency]\label{thm:binary:affine-consistency}
For fixed $p,c,\theta$,
\[
 \boxed{
 D_c(\theta)=\Delta_p(\theta)
 \ \Longleftrightarrow\
 \Tr(V_k(p,c,\theta))=0\quad\forall k\in I_\ell.
 }
\]
\end{theorem}
\begin{proof}
In characteristic two the equality
$D_c(\theta)=\Delta_p(\theta)$ is equivalent to
$D_c(\theta)+\Delta_p(\theta)=0$.
Lemma~\ref{lem:binary:V-trace} identifies
$\Tr(V_k)$ with the coordinate
$d_k(D_c+\Delta_p)$. A field element is zero exactly
when every coordinate in the basis $(\beta_k)$ is zero,
by Definition~\ref{def:binary:coords}. This proves both directions.
\end{proof}

\begin{lemma}[Mismatch gives a nonzero multilinear polynomial]\label{lem:binary:nonzero-mismatch}
If $c\neq c^p$ as tables, then
\[
 E_{c,p}(\Theta):=D_c(\Theta)+\Delta_p(\Theta)
   =\sum_{u\in I_\ell}(c_u+c^p_u)\lambda_u(\Theta)
\]
is a nonzero polynomial in $L[\Theta_0,\ldots,\Theta_{\ell-1}]$
of total degree at most $\ell$.
\end{lemma}
\begin{proof}
Each $\lambda_u$ is a product of $\ell$ affine factors,
giving the degree bound. Choose $u_0$ with
$c_{u_0}\neq c^p_{u_0}$.
Evaluate at the Boolean point $\Theta=u_0$. By
Lemma~\ref{lem:intro:interpolation} and
Lemma~\ref{lem:binary:eq-weights},
$E_{c,p}(u_0)=c_{u_0}+c^p_{u_0}\neq0$.
Hence the polynomial is nonzero.
\end{proof}

\begin{corollary}[Random-challenge detection of a coefficient mismatch]\label{cor:binary:mismatch-bound}
If $c\neq c^p$ is fixed before uniform
$\theta\in L^\ell$ is sampled, then
\[
 \Pr_\theta[D_c(\theta)=\Delta_p(\theta)]
 \leq \min\{1,\ell/q\},\qquad q=2^n.
\]
\end{corollary}
\begin{proof}
By Lemma~\ref{lem:binary:nonzero-mismatch},
$E_{c,p}$ is nonzero of total degree at most $\ell$.
The event on the left is $E_{c,p}(\theta)=0$.
Apply the multivariate Schwartz--Zippel bound proved in
Lemma~\ref{lem:history:schwartz-zippel}.
The fixed-before-challenge hypothesis ensures the
polynomial is independent of the uniform $\theta$.
\end{proof}

\begin{remark}[Meaning of this part of the construction]
The coefficients $H_{v,k}(\theta)$ turn the
$\F_2$-linear coordinate relation into \emph{affine equations
over the extension field}. This is an algebraic ingredient
of our proposed kernel-coordinate checker; it does not
by itself establish hiding of a committed witness or
the soundness of the complete ring-switching protocol.
\end{remark}

\subsection{Public trace-zero messages and their algebraic residuals}\label{sec:binary:masked-traces}
Fix $p,c,\theta$ as in Section~\ref{sec:binary:coordinates}.
The trace-zero mask of Frigo--shelat~\cite{FrigoShelat2026}
is a relevant prior algebraic mechanism; the following results
make its conditional distributions explicit for our particular
kernel-correction coefficients.

\begin{definition}[Artin--Schreier-masked public values]\label{def:binary:masked-values}
For a mask vector $t=(t_k)_{k\in I_\ell}\in L^{I_\ell}$
define
\[
 U_k(p,c,\theta,t)
   :=V_k(p,c,\theta)+t_k^2+t_k.
\]
When $t$ is random, its coordinates are sampled independently
and uniformly from $L$ unless specified otherwise.
\end{definition}

\begin{lemma}[Masking preserves the trace of each residual]\label{lem:binary:mask-trace}
For every $k\in I_\ell$,
\[
 \Tr(U_k)=\Tr(V_k)
  =d_k(D_c(\theta)+\Delta_p(\theta)).
\]
\end{lemma}
\begin{proof}
By trace additivity and
Lemma~\ref{lem:binary:as-trace},
$\Tr(U_k)=\Tr(V_k)+\Tr(t_k^2+t_k)=\Tr(V_k)$.
Apply Lemma~\ref{lem:binary:V-trace} to obtain the
second equality.
\end{proof}

\begin{corollary}[Consistency iff all public traces vanish]\label{cor:binary:masked-traces-consistent}
For fixed $p,c,\theta,t$,
\[
 D_c(\theta)=\Delta_p(\theta)
 \quad\Longleftrightarrow\quad
 \Tr(U_k)=0\quad\forall k\in I_\ell.
\]
\end{corollary}
\begin{proof}
Combine Lemma~\ref{lem:binary:mask-trace}
with Theorem~\ref{thm:binary:affine-consistency}.
\end{proof}

\begin{theorem}[Joint uniformity of independently masked traces]\label{thm:binary:masked-vector-uniform}
Suppose $D_c(\theta)=\Delta_p(\theta)$.
If the coordinates $t_k$ are independent uniform samples in $L$,
then the vector $(U_k)_{k\in I_\ell}$
is uniform on $(\ker\Tr)^{I_\ell}$.
Its distribution does not depend on the fixed values $p,c,\theta$
once the consistency equality holds.
\end{theorem}
\begin{proof}
By Theorem~\ref{thm:binary:affine-consistency},
every $V_k$ has trace zero. The random variable
$U_k=V_k+\wp(t_k)$ is therefore uniform on $\ker\Tr$
by Theorem~\ref{thm:binary:single-mask}.
The independent masks $t_k$ are transformed separately
to produce each $U_k$; hence the $U_k$ are independent.
Their joint law is the product of the uniform distributions,
which is the uniform distribution on the Cartesian product
$(\ker\Tr)^{I_\ell}$.
This product law is independent of the fixed shifts $V_k$,
and therefore of fixed $p,c,\theta$ satisfying consistency.
\end{proof}

\begin{definition}[Algebraic consistency residual]\label{def:binary:residual}
For arbitrary public values $U=(U_k)_{k\in I_\ell}\in L^{I_\ell}$
and an arbitrary mask vector $t$, set
\[
 e_k:=U_k+V_k(p,c,\theta)+t_k^2+t_k.
\]
Thus the honest masked value of Definition~\ref{def:binary:masked-values}
satisfies $e_k=0$.
\end{definition}

\begin{lemma}[Trace of an algebraic residual]\label{lem:binary:residual-trace}
For every $k$,
\[
 \Tr(e_k)=\Tr(U_k)+d_k(D_c(\theta)+\Delta_p(\theta)).
\]
\end{lemma}
\begin{proof}
By trace additivity,
$\Tr(e_k)=\Tr(U_k)+\Tr(V_k)+\Tr(\wp(t_k))$.
The last summand is zero by Lemma~\ref{lem:binary:as-trace}
and the middle summand is
$d_k(D_c+\Delta_p)$ by
Lemma~\ref{lem:binary:V-trace}.
\end{proof}

\begin{lemma}[Trace checks force a nonzero residual under mismatch]\label{lem:binary:forced-residual}
If $D_c(\theta)\neq\Delta_p(\theta)$ and
$\Tr(U_k)=0$ for all $k$, then
$e_k\neq0$ for at least one $k$.
\end{lemma}
\begin{proof}
Since $D_c(\theta)+\Delta_p(\theta)\neq0$,
at least one binary coordinate
$d_k(D_c(\theta)+\Delta_p(\theta))$ equals $1$,
by Definition~\ref{def:binary:coords}.
Lemma~\ref{lem:binary:residual-trace} and the
trace-zero assumptions give $\Tr(e_k)=1$ for
this $k$. As $\Tr(0)=0$, necessarily $e_k\neq0$.
\end{proof}

\begin{remark}[Conditional, not full-transcript hiding]
Theorem~\ref{thm:binary:masked-vector-uniform}
asserts uniformity of the public vector $U$ for
a \emph{fixed} consistent instance, assuming
independent uniform masks. It does not prove
that a later public linear combination of the
same masks $t_k$, combined with $U$, remains
witness-independent. Such correlations must
be checked in the full protocol simulation game
defined in Chapter~3.
\end{remark}

\subsection{The unanchored dependence and the anchored kernel family}\label{sec:binary:anchors}

The identity $\sum_bK_b=1$ creates an exact linear
dependence between the constant polynomial and the
kernel basis. This makes a naive residual polynomial
identity test unsound unless the additional scalar
residual is distinguished algebraically.
We derive the minimal-degree anchoring mechanism;
a full protocol using it will appear later.

\begin{lemma}[Unanchored residual dependence]\label{lem:binary:unanchored}
Let $\delta\in L$ be nonzero and set $e_b:=\delta$ for
every $b\in I_\ell$. Then
\[
 \delta+\sum_{b\in I_\ell}e_bK_b(X)=0
\]
as a polynomial, even though $(\delta,e)\neq0$.
\end{lemma}
\begin{proof}
By Lemma~\ref{lem:binary:partition},
$\sum_bK_b=1$. Therefore the left side is
$\delta+\delta\sum_bK_b=\delta+\delta=0$
in characteristic two. The first residual component
$\delta$ is nonzero by assumption.
\end{proof}

\begin{definition}[Anchored Boolean kernels]\label{def:binary:anchored}
Let $0_\ell=(0,\ldots,0)\in I_\ell$. For $b\in I_\ell$ define
\[
 W_b(X):=
 \begin{cases}
 K_{0_\ell}(X)+X^n,&b=0_\ell,\\
 K_b(X),&b\neq0_\ell.
 \end{cases}
\]
Every $W_b$ is a polynomial with coefficients in $\F_2$.
\end{definition}

\begin{lemma}[Anchored family degree]\label{lem:binary:anchor-degree}
$\deg W_{0_\ell}=n$, while
$\deg W_b=n-1$ for every $b\neq0_\ell$.
\end{lemma}
\begin{proof}
By Lemma~\ref{lem:intro:degree}, each kernel has degree
$n-1$. Therefore the added monomial $X^n$ is the
unique degree-$n$ contribution to $W_{0_\ell}$,
with nonzero coefficient $1$, so its degree is $n$.
For $b\neq0_\ell$, $W_b=K_b$, which has degree $n-1$.
\end{proof}

\begin{lemma}[Sum of anchored kernels]\label{lem:binary:anchor-sum}
\[
 \sum_{b\in I_\ell}W_b(X)=1+X^n.
\]
\end{lemma}
\begin{proof}
By definition only the term $W_{0_\ell}$
differs from $K_{0_\ell}$, by the summand $X^n$.
Thus $\sum_bW_b=\sum_bK_b+X^n=1+X^n$
by Lemma~\ref{lem:binary:partition}.
\end{proof}

\begin{theorem}[Independence of the constant and all anchored kernels]\label{thm:binary:anchor-independence}
The $n+1$ polynomials
\[
 \boxed{\{1\}\cup\{W_b:b\in I_\ell\}}
\]
are linearly independent over $L$ in $L[X]_{\leq n}$.
\end{theorem}
\begin{proof}
Suppose
\[
 \delta+\sum_{b\in I_\ell}e_bW_b(X)=0
 \qquad(\delta,e_b\in L).
\]
The coefficient of $X^n$ in the left side is
exactly $e_{0_\ell}$ by Lemma~\ref{lem:binary:anchor-degree}.
Hence $e_{0_\ell}=0$. The remaining polynomial identity is
$\delta+\sum_{b\neq0_\ell}e_bK_b=0$.
Using $1=\sum_bK_b$ from
Lemma~\ref{lem:binary:partition}, rewrite it as
\[
 \delta K_{0_\ell}
 +\sum_{b\neq0_\ell}(\delta+e_b)K_b=0.
\]
By the kernel basis theorem,
Proposition~\ref{prop:intro:basis},
all kernel coefficients vanish. In particular $\delta=0$,
and then $\delta+e_b=0$ gives every $e_b=0$.
Thus the anchored family together with $1$ is independent.
\end{proof}

\begin{corollary}[Anchored family is a basis in degree at most $n$]\label{cor:binary:anchor-basis}
The $n+1$ polynomials from
Theorem~\ref{thm:binary:anchor-independence}
form an $L$-basis of $L[X]_{\leq n}$.
\end{corollary}
\begin{proof}
The monomials $1,X,\ldots,X^n$ form a basis of
$L[X]_{\leq n}$, which therefore has dimension $n+1$.
The displayed family consists of $n+1$ linearly
independent vectors in this space, hence is a basis.
\end{proof}

\begin{proposition}[Degree-$n$ anchoring is dimensionally minimal]\label{prop:binary:anchor-minimal}
No family of $n+1$ polynomials, all of degree $<n$,
can be linearly independent over $L$.
\end{proposition}
\begin{proof}
The space $L[X]_{<n}$ has basis
$1,X,\ldots,X^{n-1}$ and therefore dimension $n$.
A vector space of dimension $n$ cannot contain $n+1$
linearly independent vectors.
\end{proof}

\begin{lemma}[Anchored kernels commute with Frobenius evaluation]\label{lem:binary:anchor-frob}
For every $b\in I_\ell$ and $z\in L$,
\[
 W_b(z)^2=W_b(z^2).
\]
\end{lemma}
\begin{proof}
Every coefficient of $W_b$ lies in $\F_2$
by Definition~\ref{def:binary:anchored}.
The coefficient Frobenius of $W_b$ is therefore $W_b$.
Apply Lemma~\ref{lem:binary:poly-frob} with $j=1$
and evaluate at $z$:
$W_b(z)^2=W_b^{\sigma_1}(z^2)=W_b(z^2)$.
\end{proof}

\begin{definition}[Anchored residual polynomial]\label{def:binary:anchor-residual}
For $\delta\in L$ and $e=(e_b)_{b\in I_\ell}\in L^{I_\ell}$,
define
\[
 F_{\delta,e}(X):=\delta+\sum_{b\in I_\ell}e_bW_b(X).
\]
\end{definition}

\begin{lemma}[Nonzero residual gives a nonzero polynomial]\label{lem:binary:anchor-nonzero}
If $(\delta,e)\neq(0,0)$, then
$F_{\delta,e}$ is a nonzero polynomial of degree at most $n$.
\end{lemma}
\begin{proof}
Each $W_b$ has degree at most $n$,
so the linear combination has degree at most $n$.
If it vanished identically, the linear independence
of $\{1,W_b\}_b$ proved in
Theorem~\ref{thm:binary:anchor-independence}
would force $\delta=0$ and $e_b=0$ for all $b$,
contradicting the hypothesis.
\end{proof}

\begin{theorem}[Root bound for an anchored residual]\label{thm:binary:anchor-root}
Fix $(\delta,e)\neq(0,0)$ independently of a uniform
challenge $z\overset{\$}{\leftarrow}L$. Then
\[
 \boxed{\Pr_z\left[
 \delta+\sum_be_bW_b(z)^2=0\right]\leq\frac{n}{q}.}
\]
\end{theorem}
\begin{proof}
By Lemma~\ref{lem:binary:anchor-frob},
the tested expression is $F_{\delta,e}(z^2)$.
By Corollary~\ref{cor:binary:frob-bijective},
the map $z\mapsto z^2$ is a bijection of the finite
field $L$, so $z^2$ is uniform on $L$.
By Lemma~\ref{lem:binary:anchor-nonzero},
$F_{\delta,e}$ is nonzero of degree at most $n$.
The univariate root bound of
Lemma~\ref{lem:history:rootcount} gives at most $n$
roots. Divide by $q$.
\end{proof}

\begin{remark}[Challenge timing is necessary]
The bound is conditional on $(\delta,e)$ being
fixed independently of the fresh challenge $z$.
If the coefficients are chosen after seeing $z$,
the prover could select a nonzero polynomial
that vanishes at that particular point.
Thus the commit-then-challenge order is part
of the soundness hypothesis, not a cosmetic
presentation detail.
\end{remark}

\subsection{An algebraic two-challenge consistency checker}\label{sec:binary:checker}
We now combine the preceding lemmas into a small
\emph{algebraic checking theorem}. This is the
kernel-coordinate consistency component of the
proposed binary-field construction. The theorem
does not assert zero knowledge for the complete IOP
or its commitment scheme.

\begin{definition}[Affine consistency checks]\label{def:binary:checks}
Let the secret values $p,c,t$, $r,s$ have types
\[
 p,c,t\in L^{I_\ell},\qquad r,s\in L.
\]
After a challenge $\theta\in L^\ell$,
let $V_k=V_k(p,c,\theta)$ be as in
Definition~\ref{def:binary:affine-H}, and let
$U=(U_k)_k\in L^{I_\ell}$ be public.
After a second challenge $z\in L$, a public scalar
$A\in L$ is supplied. The \emph{algebraic
consistency checker} accepts precisely if
\begin{align}
 \Tr(U_k)&=0\quad\text{for all }k\in I_\ell,
 \label{eq:binary:check-traces}\\
 A+r+\sum_kW_k(z)t_k&=0,
 \label{eq:binary:check-first}\\
 A^2+s+\sum_kW_k(z)^2(U_k+V_k+t_k)&=0.
 \label{eq:binary:check-second}
\end{align}
The last two equations are affine in the
committed secret variables $p,c,t,r,s$ once
$\theta,z,U,A$ are fixed: $A^2$ and $W_k(z)^2$
are public scalars.
\end{definition}

\begin{definition}[Scalar and vector checking residuals]\label{def:binary:check-residuals}
Put
\[
 \delta:=s+r^2,\qquad
 e_k:=U_k+V_k+t_k^2+t_k
\]
and form the anchored residual polynomial
$F_{\delta,e}$ of Definition~\ref{def:binary:anchor-residual}.
\end{definition}

\begin{lemma}[First equation determines the public scalar]\label{lem:binary:first-determines}
For fixed $r,t,z$, equation~\eqref{eq:binary:check-first}
holds if and only if
\[
 A=r+\sum_kW_k(z)t_k.
\]
\end{lemma}
\begin{proof}
Move the other terms to the right side of
equation~\eqref{eq:binary:check-first}.
In characteristic two, additive inverses
equal the original values, so subtraction and
addition coincide. The resulting formula is
both necessary and sufficient.
\end{proof}

\begin{lemma}[The two affine checks imply an anchored residual root]\label{lem:binary:check-implies-root}
If equations~\eqref{eq:binary:check-first}
and~\eqref{eq:binary:check-second} both hold, then
\[
 \boxed{\delta+\sum_ke_kW_k(z)^2=0.}
\]
\end{lemma}
\begin{proof}
By Lemma~\ref{lem:binary:first-determines},
$A=r+\sum_kW_k(z)t_k$. Square this identity.
Characteristic-two additivity of squares
(Lemma~\ref{lem:binary:square}) gives
\[
 A^2=r^2+\sum_kW_k(z)^2t_k^2.
\]
Substitute it into equation~\eqref{eq:binary:check-second}:
\[
 0=r^2+s+
   \sum_kW_k(z)^2\bigl(t_k^2+U_k+V_k+t_k\bigr).
\]
The first two terms are $\delta=s+r^2$,
and the parenthesis is exactly $e_k$ by
Definition~\ref{def:binary:check-residuals}.
This is the claimed identity.
\end{proof}

\begin{lemma}[Honest trace check]\label{lem:binary:honest-traces}
Assume $c=c^p$ and choose
$U_k:=V_k+\wp(t_k)$.
Then the trace equations~\eqref{eq:binary:check-traces}
hold for every $\theta$ and every $t$.
\end{lemma}
\begin{proof}
By the definition of $c^p$ and
Definition~\ref{def:binary:transpose},
$D_c(\theta)=\Delta_p(\theta)$ for every $\theta$.
Corollary~\ref{cor:binary:masked-traces-consistent}
then yields $\Tr(U_k)=0$ for all $k$.
\end{proof}

\begin{lemma}[Honest affine checks]\label{lem:binary:honest-affine}
Assume $c=c^p$, $s=r^2$, and
$U_k=V_k+\wp(t_k)$ for every $k$.
For every $\theta,z$, choose
\[
 A:=r+\sum_kW_k(z)t_k.
\]
Then the two affine equations
\eqref{eq:binary:check-first} and
\eqref{eq:binary:check-second} hold.
\end{lemma}
\begin{proof}
The first equation holds by definition of $A$,
since $A+A=0$ in characteristic two.
The second equation follows from
$U_k+V_k+t_k=t_k^2$:
\[
 A^2+s+\sum_kW_k(z)^2(U_k+V_k+t_k)
 =A^2+r^2+\sum_kW_k(z)^2t_k^2.
\]
By the same squaring calculation used in
Lemma~\ref{lem:binary:check-implies-root},
$A^2=r^2+\sum_kW_k(z)^2t_k^2$.
The displayed sum therefore consists of two
identical copies of each term and equals zero.
\end{proof}

\begin{theorem}[Perfect algebraic completeness]\label{thm:binary:checker-complete}
If $c=c^p$, $s=r^2$, $U_k=V_k+\wp(t_k)$, and
$A=r+\sum_kW_k(z)t_k$, the checker accepts for
every choice of the challenges $(\theta,z)$.
\end{theorem}
\begin{proof}
Lemma~\ref{lem:binary:honest-traces} proves the
trace tests, and Lemma~\ref{lem:binary:honest-affine}
proves the two remaining affine equations.
Thus all tests in Definition~\ref{def:binary:checks}
hold.
\end{proof}

\begin{theorem}[Conditional anchored-check soundness]\label{thm:binary:checker-conditional}
Fix $\theta,p,c,t,r,s,U$ before a uniform
$z\overset{\$}{\leftarrow}L$ is sampled.
Suppose the trace checks hold and
$D_c(\theta)\neq\Delta_p(\theta)$.
Then, even if $A$ is chosen after seeing $z$,
the probability that both affine checks pass
is at most $n/q$.
\end{theorem}
\begin{proof}
Because the traces of all $U_k$ vanish while
$D_c(\theta)\neq\Delta_p(\theta)$,
Lemma~\ref{lem:binary:forced-residual}
shows that some $e_k$ is nonzero.
Thus $(\delta,e)\neq(0,0)$.
The residual values depend only on
$\theta,p,c,t,r,s,U$, which are fixed before $z$,
so the hypotheses of Theorem~\ref{thm:binary:anchor-root}
apply. If both affine checks pass,
Lemma~\ref{lem:binary:check-implies-root}
forces
$\delta+\sum_ke_kW_k(z)^2=0$.
The anchored-root bound gives probability
at most $n/q$, independently of how $A$
is selected after the challenge.
\end{proof}

\begin{theorem}[Two-challenge coefficient-consistency bound]\label{thm:binary:two-challenge}
Fix $p,c\in L^{I_\ell}$ with $c\neq c^p$ before
sampling uniform $\theta\in L^\ell$.
After seeing $\theta$, an arbitrary strategy
may select $t,r,s,U$ (and must satisfy all trace
checks), but these values must be fixed
\emph{before} an independent uniform
$z\overset{\$}{\leftarrow}L$ is sampled.
The strategy may choose $A$ after seeing $z$.
Then the probability that all checks in
Definition~\ref{def:binary:checks} accept is at most
\[
 \boxed{
 \min\left\{1,\frac{\ell+n}{q}\right\},
 \qquad q=2^n.
 }
\]
\end{theorem}
\begin{proof}
Let $E$ be the event
$D_c(\theta)=\Delta_p(\theta)$.
Since $c\neq c^p$ was fixed before $\theta$,
Corollary~\ref{cor:binary:mismatch-bound}
gives $\Pr[E]\leq\ell/q$.
On the complementary event $E^c$,
condition on any particular $\theta$ and on
the strategy's chosen values $t,r,s,U$
(and its private random tape), provided
the public trace checks pass.
These values are fixed before the fresh
challenge $z$, so
Theorem~\ref{thm:binary:checker-conditional}
bounds the conditional probability of
passing both affine checks by $n/q$.
If the trace checks fail, the acceptance
probability is zero. Averaging over these
conditional cases yields
\[
 \Pr[\text{accept}]
 \leq\Pr[E]+\Pr[\text{accept}\cap E^c]
 \leq\frac{\ell}{q}+\frac{n}{q}.
\]
Probabilities are also bounded by $1$.
This proves the stated minimum.
\end{proof}

\begin{lemma}[No even-order multiplicative subgroup in a binary field]\label{lem:binary:no-even-subgroup}
Every subgroup of $L^\times$ has odd order.
In particular, the odd-characteristic
pairing $x\leftrightarrow-x$ used in the
original KBFold word fold has no distinct
paired points in $L^\times$.
\end{lemma}
\begin{proof}
The group $L^\times$ has order
$|L|-1=2^n-1$, an odd integer.
The order of any subgroup divides
the order of the finite ambient group,
so every subgroup has odd order.
Furthermore $-x=x$ for every $x\in L$
because the characteristic is two.
Thus $x$ and $-x$ cannot form two distinct
points.
\end{proof}

\begin{remark}[Precise boundary of the result]
Theorems~\ref{thm:binary:checker-complete} and
\ref{thm:binary:two-challenge} concern only
the \emph{algebraic coefficient-consistency
checker}. They assume the stated timing
and do not incorporate a concrete
commitment-binding error, the remaining
ring-switching equations, the simulator
for the joint public messages $(U,A)$,
or any proof of security for the
odd-characteristic KBFold IOP.
The final lemma also proves that the
binary checker cannot simply reuse the
$\pm x$ multiplicative-domain fold.
These issues will be addressed in the
appropriate later chapters.
\end{remark}


\subsection{The theorem chain and its formal boundary}\label{sec:binary:conclusion}

For $L=\F_{2^n}$ with $n=2^\ell$ and
$c,p\in L^{I_\ell}$, the algebraic chain is
\[
 \boxed{
 \Delta_p(\theta)
 =D_{C_p}(\theta)
 =\mathsf{PFold}^{[\ell]}_\theta(C_p)
 =\mathscr M_\ell(C_p)(1+\theta).
 }
\]
The correction can be verified through traces of affine
expressions over $L$:
\[
 \boxed{
 D_c(\theta)=\Delta_p(\theta)
 \iff
 \Tr(V_k(p,c,\theta))=0\quad\forall k.
 }
\]
The Artin--Schreier map $\wp(t)=t^2+t$ has image
$\ker\Tr$, and independent uniform masks make the
\emph{individual public trace vector} uniformly distributed
over $(\ker\Tr)^{I_\ell}$ for a fixed consistent instance.
The anchored family
\[
 W_b=K_b+\mathbf 1_{\{b=0_\ell\}}X^n
\]
satisfies that $\{1,W_b:b\in I_\ell\}$ is a basis of
$L[X]_{\leq n}$. With secret coefficients and messages
fixed before the relevant fresh challenges, the
two-challenge consistency checker has mismatch-acceptance
probability at most
\[
 \boxed{\min\{1,(\ell+n)/2^n\}.}
\]
This bound concerns algebraic coefficient consistency alone.
The complete security claim requires the remaining
ring-switching equations, binding commitments, a
joint-transcript simulator, and the composition analysis
developed in later chapters. No new Lean certificate is
asserted here; each result is mapped to a proposed Lean
declaration in the chapter ledger.
